\section{Relation to prior work}\label{sec:literature}

\subsection{Prescribed dg kernels}
The construction in \Cref{thm:realisation} belongs to the established setting of matrix-ring universal localisation. Sheiham \cite[Theorem~2.4 and Example~2.5(1)]{Sheiham06} gives explicit localisations of triangular matrix rings at a map between column projectives, including the passage from an upper triangular ring to a full matrix ring. Morita contexts and exact-context localisations are studied in \cite{GreenPsaroudakis14,GaoPsaroudakis17,ChenXiRecollI,ChenXi19}. The claim here is the prescribed identification of the kernel dg algebra, by a strict dg map, together with two-sided target projectivity, t-exactness, and full profile. The matrix-ring mechanism by itself does not assert that package.

Marks--Vit\'oria \cite[Theorem~3.3]{MarksVitoria15} prove that a ring epimorphism whose target is finitely presented as a left source module of projective dimension at most one is homological if and only if it is a universal localisation. In particular, their recognition theorem applies to the projective targets constructed here. Our t-exact converse instead imposes a condition on a marked compact kernel generator and forces that generator to be two-term; it does not assume a projective-dimension bound on the target. Its universal-localisation conclusion should be read alongside this earlier recognition result.

Minimal \(A_\infty\)-models and their uniqueness are standard \cite{Keller01Ainf,LefevreHasegawa03}; crossed extensions are classified by \(HH^3\) in \cite{BauesMinian02}. Finite strict models for proper connective dg algebras over arbitrary fields are supplied by Raedschelders--Stevenson's appendix to Goodbody \cite[Theorem~A.3]{Goodbody24}. Accordingly, \Cref{thm:finite-crossed} is an explicit finite construction from a Hochschild class, rather than a new finite-model existence theorem. The ordinary-ring kernel realisation, including its non-formal example, is the separate assertion.

Kalck--Yang \cite{KalckYangI,KalckYangII} and Booth \cite{Booth21} describe derived quotients associated with idempotents. Their Schur-functor recollements are a different realisation setting from extension along an ordinary ring epimorphism. The compact projection in \Cref{prop:left-projection} and the negative Calabi--Yau argument isolate a boundary between these settings; \Cref{ex:sphere-schur} makes it explicit. No general converse for Schur kernels is proved here.

\subsection{Detectors and changing profiles}
The finite-total-Hom characterisation is recorded in \cite[Lemma~4.1]{GaoPsaroudakis18}; intrinsic homological finiteness is further developed in \cite{KuznetsovShinder25}. Far-away orthogonality and Brown--Comenetz detection provide the framework of \cite{KostasPsaroudakisVitoria25}. The canonical detector in \Cref{thm:detector} applies the familiar duality \(\RHom_R(X,DR)\simeq D_kX\) to the embedded kernel. Its purpose is to identify the additional data determining ambient boundedness.

One-point extensions are classical \cite{Ringel84}, and lifting recollements along them is studied in \cite{LinLin08}. \Cref{thm:ambient} computes the extended universal localisation and its unchanged marked kernel, proves the all-degree stable-flatness criterion, and identifies the new profile as an intersection of detector cuts. The degree-zero concentration test comes from \cite{NicolasSaorin09}. Chen's triangular singularity equivalence \cite[Theorem~4.1(1)]{ChenSchur09} supplies the ambient singularity category for the ordinary-kernel examples. Marks' Nakayama classification \cite{Marks15} and the relation of universal localisation to silting \cite{MarksStovicek19} are used with the stated module-side and field hypotheses.

\subsection{Classical comparison and its refinements}
The triangulated compact-kernel comparison is a classical consequence, not a new theorem requiring dg methods. Krause's compact factorisation theorem \cite[Theorem~7.2.1(1)]{Krause10} supplies the source-compact orientation of his Lemma~4.7.1(1), giving \Cref{thm:abstract-rigidity} without algebraicity. The opposite orientation of the quotient criterion is also standard \cite{Keller96,Krause10,SunZhangZhang}. Efimov's Calkin exactness \cite[Proposition~1.4]{Efimov17} is the enhanced mechanism retained in \Cref{sec:obstruction}.

Singularity localisation in recollements is developed in \cite{JinYangZhou,SunZhangZhang}. The bounded hypotheses in \Cref{thm:actual-kernel} identify the displayed ambient quotient with the actual object-kernel of a singularity functor. Schlichting localisation, the completeness criteria discussed in \cite{KalckThesis,PavicShinder21}, and connective invariance \cite{BGT13} are the inputs to \Cref{app:negative-k}. The simultaneous kernel formula compares two nested dense quotients, rather than supplying a new negative-\(K\) localisation sequence. Its obstruction vanishes for all examples in this paper.

\subsection{Support applications}
The support intersection criterion and realisability over complete intersections are from \cite{AvramovBuchweitz00,BerghJorgensenEasy}; thick subcategories are classified by support in \cite{Stevenson14}. Hyperplane avoidance is identified with matrix factorisations in \cite[Theorem~3.5]{BerghJorgensenMF}; higher-codimension constructions appear in \cite{BurkeWalker15}. For \(p\)-groups the inputs include \cite{Carlson83,Carlson84,AvruninScott82,Quillen71}. The application here is the realisation of these marked subcategories as bounded defects of ordinary finite-dimensional algebras, together with the no-generator consequence. It does not classify arbitrary abstract equivalences between the resulting categories, as \Cref{ex:support-symmetry} illustrates.
